% Motivation letter — University of Bologna, Bachelor's in Digital Humanities and Digital Knowledge

% Shared LaTeX preamble for all motivation letters.
% Don't compile this file directly --- it's imported by each per-university .tex.
\documentclass[11pt,a4paper]{letter}

\usepackage[a4paper,margin=2.2cm]{geometry}
\usepackage[T1]{fontenc}
\usepackage[utf8]{inputenc}
\usepackage[sfdefault]{roboto}
\usepackage[hidelinks]{hyperref}
\usepackage{xcolor}
\usepackage{microtype}
\usepackage{parskip}

\definecolor{accent}{HTML}{0e7490}

\hypersetup{
  colorlinks=true,
  urlcolor=accent,
}

% Legal name used here — required on university applications.
% Professional name "Ahmad Bagheri" is used on LinkedIn / resume / OSS.
\name{Ahmad Baghereslami}
\address{
  contact@ahmadcodes.com \\
  +37494308785 \\
  ahmadcodes.com \\
  Armenia
}
\signature{Ahmad Baghereslami}

\input{glyphtounicode}
\pdfgentounicode=1


\begin{document}

\begin{letter}{
  International Admissions Office\\
  Alma Mater Studiorum --- Università di Bologna\\
  Via Zamboni 33\\
  40126 Bologna, Italy
}

\opening{Dear Admissions Committee,}

I am applying for the Bachelor's in Digital Humanities and Digital Knowledge at the University of Bologna for the September 2027 intake. This programme is an unusually precise fit for my background, and this letter explains why.

I trained in the Humanities track (Olum-e-Ensani) at high school in Iran and earned my Pre-University Certificate in 2014. I did not sit for the Konkur examination; I began programming professionally at eighteen instead. Over the seven years since, my career has been an unintentional bridge between two worlds: humanistic training in language, writing, and structured argument on one side, and applied software engineering on the other. Most days the latter dominates, but the way I approach problems --- mapping data flow as if it were narrative, paying attention to the meaning of an API's contract as much as its mechanics --- carries the imprint of the former.

The work I am most proud of reflects this. At Barriertek, a fire-retardant lumber producer in Toronto, I am the sole engineer on Woody Portal, a customer portal that replaces phone-tag, whiteboard production scheduling, and per-team Excel sheets across distributors, truckers, production workers, and the office team. End-to-end TypeScript, bi-directional NetSuite integration via SuiteScript RESTlets, deployed on PostgreSQL and Linux. The technical surface is conventional. What is less conventional is how much of my time goes into translating between the formal data model of an enterprise ERP and the informal vocabulary of a 50-year-old manufacturer --- a translation problem that is essentially humanistic. I open-sourced the integration client from this project at \href{https://github.com/Ahmad-b1995/netsuite-restlet-typescript}{github.com/Ahmad-b1995/netsuite-restlet-typescript}.

Bologna's Digital Humanities and Digital Knowledge programme is, on paper, the formalization of the arc my work has already traced. The integration of computational thinking, linguistic processing, and the methodologies of cultural data is rare at the undergraduate level. The Department of Classical Philology and Italian Studies' partnership with Computer Science is exactly the framework I want to study in. I have been reading the syllabus, the work of professors Fabio Tamburini and Francesca Tomasi, and the public output of the /DH.arc lab; the questions you take seriously --- about how the technical and the textual constrain each other --- are the questions my professional work has stumbled into without the right vocabulary.

I am also drawn to Bologna for the obvious reasons. The Alma Mater Studiorum's age and the seriousness of the university's intellectual culture matter. So does Bologna as a city: a student city with affordable living relative to Milan, a strong ER.GO scholarship system, and an international cohort large enough that an Iranian student arriving with a partner can find footing without isolation.

I will spend my student visa's twenty weekly hours of work on contract software engineering in my specialty (ERP integration, NetSuite), keeping my professional practice alive alongside coursework. After graduation I plan to remain in the EU, ideally building a small consultancy in software-mediated communication for traditional industrial businesses --- the same kind of work I am doing now, but with the formal grounding to do it better.

I would be grateful for the opportunity to study at Bologna.

\closing{Sincerely,}

\end{letter}

\end{document}
